% Wygląd wkładki: akapit, kolory, etykiety, pagina, okładka, spis utworów.
\makeatletter

% ---------- Akapit -----------------------------------------------------
\linespread{1.06}
\setlength{\parindent}{1.1em}
\frenchspacing
\clubpenalty=10000
\widowpenalty=10000
\displaywidowpenalty=10000

% ---------- Kolory -----------------------------------------------------
\definecolor{bxpapier}{HTML}{F3EBD7}
\definecolor{bxatrament}{HTML}{2A231F}
\definecolor{bxszary}{HTML}{8A7F72}
% paleta tierów
\definecolor{musztarda}{HTML}{D39B22}
\definecolor{pomarancz}{HTML}{C9602A}
\definecolor{rdza}{HTML}{973A25}
\definecolor{sliwka}{HTML}{5A3354}
\definecolor{morski}{HTML}{2D6B67}
\definecolor{oliwka}{HTML}{6C7733}
\colorlet{tier}{rdza}

\ifekran \pagecolor{bxpapier}\fi
\AtBeginDocument{\color{bxatrament}}

% ---------- Etykiety ---------------------------------------------------
\ifpolski
\def\bx@lside{Strona}\def\bx@ltracklist{Spis utworów}
\else
\def\bx@lside{Side}\def\bx@ltracklist{Tracklist}
\fi

% ---------- Dane okładki -----------------------------------------------
\def\bx@tytul{Tytuł}\def\bx@podtytul{}\def\bx@autor{}\def\bx@stopka{}
\def\bx@tytulkrotki{}
\newcommand\tytul[2][]{\def\bx@tytul{#2}\def\bx@tytulkrotki{#1}}
\newcommand\podtytul[1]{\def\bx@podtytul{#1}}
\newcommand\autor[1]{\def\bx@autor{#1}}
\newcommand\stopka[1]{\def\bx@stopka{#1}}

% tytuł okładki: 40 pt, ale zmniejszany, gdy najdłuższy wiersz
% nie mieści się między marginesami okładki (14 mm z każdej strony)
\newsavebox\bx@tytulbox
\newcommand\bx@tytulfit{%
  \sbox\bx@tytulbox{\bxdisplay\fontsize{40}{37}\selectfont
    \begin{tabular}[t]{@{}l@{}}\bx@tytul\end{tabular}}%
  \ifdim\wd\bx@tytulbox>\dimexpr\paperwidth-28mm\relax
    \resizebox{\dimexpr\paperwidth-28mm\relax}{!}{\usebox\bx@tytulbox}%
  \else\usebox\bx@tytulbox\fi}

% rozstrzelone wersaliki w kroju groteskowym
\newcommand\bxcaps[2][14]{{\sffamily\addfontfeature{LetterSpace=#1}\MakeUppercase{#2}}}

% pasy z lat 70.: [grubość]{lista kolorów}{odległość od góry strony}
\newcommand\bx@pasy[3][5mm]{%
    \begin{tikzpicture}[remember picture, overlay]
        \foreach \kol [count=\i] in {#2} {
            \fill[\kol] ([yshift={-#3-(\i-1)*#1}]current page.north west)
            rectangle ([yshift={-#3-\i*#1}]current page.north east);
        }
\end{tikzpicture}}

% ---------- Liczniki ---------------------------------------------------
\newcounter{bxside}
\newcounter{bxtrack}[bxside]

% ---------- Nagłówki i stopki ------------------------------------------
\def\bx@mark{}
\setkomafont{pageheadfoot}{\normalfont}
\pagestyle{scrheadings}
\clearpairofpagestyles
\lehead{\color{bxszary}\bxcaps[18]{\scriptsize\bx@tytulkrotki}}
\rohead{\color{tier}\bxcaps[18]{\scriptsize\bx@mark}}
\lefoot*{\sffamily\footnotesize\thepage}
\rofoot*{\sffamily\footnotesize\thepage}

% inicjał w stylu wkładki
\renewcommand{\LettrineFontHook}{\bxdisplay\color{tier}}
\renewcommand{\LettrineTextFont}{\sffamily\addfontfeature{LetterSpace=8}\MakeUppercase}
\setcounter{DefaultLines}{3}
\setlength{\DefaultFindent}{2pt}
\setlength{\DefaultNindent}{0pt}

% ========================================================================
\ExplSyntaxOn

% ---------- Okładka ----------------------------------------------------
\NewDocumentCommand \okladka { }
{
    \begin{titlepage}
        \bx@pasy[7mm]{musztarda,pomarancz,rdza,sliwka,morski,oliwka}{124mm}
        \begin{tikzpicture}[remember~picture, overlay]
            \node[anchor=north~west, text=bxatrament, inner~sep=0pt]
            at ([xshift=14mm, yshift=-22mm]current~page.north~west)
            {\bx@tytulfit};
            \node[anchor=north~west, align=left, text=bxatrament, inner~sep=0pt,
            font=\itshape\fontsize{12.5}{15.5}\selectfont, text~width=118mm]
            at ([xshift=14.5mm, yshift=-90mm]current~page.north~west)
            {\bx@podtytul};
            \node[anchor=south~west, align=left, text=bxatrament, inner~sep=0pt]
            at ([xshift=14.5mm, yshift=18mm]current~page.south~west)
            {\bxcaps[22]{\footnotesize\bx@autor}};
            \node[anchor=south~east, align=right, text=bxszary, inner~sep=0pt]
            at ([xshift=-14mm, yshift=18mm]current~page.south~east)
            {\bxcaps[22]{\scriptsize\bx@stopka}};
        \end{tikzpicture}
        \null
    \end{titlepage}
}

% ---------- Spis utworów (tył okładki LP), rozkładówka -----------------
\NewDocumentCommand \tracklista { }
{
    \cleardoubleevenemptypage
    \begingroup
    \pagestyle{empty}
    \noindent{\bxdisplay\fontsize{22}{24}\selectfont \bx@ltracklist}\par
    \vspace{2mm}
    \raggedright
    \use:c{@starttoc}{trk}
    \clearpage
    \endgroup
}

\NewDocumentCommand \bxtrkside { m m m m }
{
    \par\addvspace{8pt}
    \Needspace{11\baselineskip} % cała strona płyty (6 pozycji) w jednym kawałku
    \colorlet{tier}{#4}
    \noindent
    {\color{#4}\bxcaps[16]{\scriptsize\bfseries \bx@lside{}~#1}}
    \hspace{2mm}{\bxdisplay\small #2}
    \tl_if_blank:nF {#3} { \hspace{1.5mm}{\color{bxszary}\itshape\footnotesize #3} }
    \par\vspace{1pt}
    {\color{#4}\rule{\linewidth}{0.9pt}}\par\vspace{2.5pt}
}

\NewDocumentCommand \bxtrkentry { m m m m m m }
{
    \par\noindent
    \makebox[7.5mm][l]{\hyperlink{bx-#1}{\color{tier}\sffamily\bfseries\scriptsize #1}}%
    \parbox[t]{\dimexpr\linewidth-7.5mm-7mm}
    {
        \raggedright\linespread{0.95}\selectfont
        {\sffamily\small #2}
        \tl_if_blank:nF {#3} { ~{\color{bxszary}\sffamily\scriptsize (#3)} }
        \hspace{0.4em}\textcolor{bxszary}{·}\hspace{0.4em}
        {\itshape\small #4}
        \tl_if_blank:nF {#5} { ~{\color{bxszary}\itshape\scriptsize (#5)} }
    }%
    \hfill
    \makebox[7mm][r]{\hyperlink{bx-#1}{\sffamily\scriptsize #6}}
    \par\vspace{1.8pt}
}

\ExplSyntaxOff
% ========================================================================

\makeatother
